%% Anexo A: resumen de estimación de costos. Los cálculos completos
%% provienen de sistema/estimacion/ y del tablero del proyecto.
\chapter{Resumen de estimación de costos}
\label{anexo:estimacion}

Los valores siguientes son supuestos académicos de planificación expresados
en bolivianos (Bs). No se cuenta con cotizaciones reales de mercado: la
tarifa y la capacidad son las declaradas por el equipo, y el ritmo medido
proviene del tablero de trabajo.

\section{COCOMO básico}

Se supone un proyecto orgánico de 3 KLOC (3.000 líneas de código) y un costo
de 4.000 Bs por persona-mes. Se emplean las ecuaciones de COCOMO básico
\parencite{boehm1981}, con los resultados de la tabla~\ref{tab:cocomo}.

\begin{tablaapa}{tab:cocomo}{Estimación con COCOMO básico}
\small
\begin{tabular}{@{}p{4.6cm}p{5.6cm}p{4.4cm}@{}}
\toprule
Concepto & Cálculo & Resultado \\
\midrule
Esfuerzo & $2{,}4 \times 3^{1{,}05}$ & 7,61 persona-mes \\
Duración & $2{,}5 \times 7{,}61^{0{,}38}$ & 5,40 meses \\
Personal promedio & $7{,}61 \div 5{,}40$ & 1,41 personas \\
Costo laboral & $7{,}61 \times 4.000$ Bs & 30.426,26 Bs \\
\bottomrule
\end{tabular}
\nota{Elaboración propia con \textcite{boehm1981}.}
\end{tablaapa}

\section{Estimación ágil}

Se utiliza el alcance real del backlog (19 historias, 55 puntos de historia y
54 tareas) y el ritmo medido en el tablero durante la primera ventana de
trabajo: 4,00 horas-persona por punto de historia. La tarifa declarada es de
5.000 Bs por mes-persona de 160 horas; la capacidad del equipo es de 8
horas-persona por día hábil (dos integrantes con 4 horas). Los resultados se
resumen en la tabla~\ref{tab:agil}.

\begin{tablaapa}{tab:agil}{Estimación ágil con datos del proyecto}
\small
\begin{tabular}{@{}p{4.6cm}p{5.6cm}p{4.4cm}@{}}
\toprule
Concepto & Cálculo & Resultado \\
\midrule
Producto mínimo viable de la primera ventana & 8 pts (HU-08, HU-01 y
HU-04) & 2.500 Bs \\
Esfuerzo del alcance completo & $55 \times 4{,}00$ h = 220 h & 1,38
mes-persona \\
Costo del alcance completo & $1{,}38 \times 5.000$ Bs & 6.875 Bs \\
Calendario & $220 \div 8$ h por día de equipo & 27,5 días hábiles (2,75
ventanas) \\
Cierre estimado & desde el 25/11/2026 & 07/01/2027 \\
\bottomrule
\end{tabular}
\nota{Elaboración propia a partir del backlog y del tablero del proyecto.}
\end{tablaapa}

Los dos modelos no son directamente equivalentes: COCOMO emplea un esfuerzo
equivalente a tiempo completo sobre un tamaño hipotético de código, y la
estimación ágil parte del alcance real y del ritmo medido con dedicación
parcial. Se adopta el cálculo ágil como presupuesto preliminar del proyecto y
se recalcula al cierre de cada ventana de trabajo. COCOMO no incluye
infraestructura ni contingencia.
